\section{Clearance and contour estimates}
\label{app:estimates}

We give explicit constants for the spatial and index contours in
Sections~\ref{sec:complex} and \ref{sec:cauchy}. The source inputs
are \eqref{eq:domain} and \eqref{eq:count-safe}
\cite{BakerBase,BakerGT41}. Write \(q_0=16/15\) and
\(\Lambda(x)=\log(3+2T+|x|)\).

\subsection{Complete value sums}\label{app:value-sums}

Let \(p=u+\ii v\), \(u>1/2\). The band \(|\Re w_j|<1\) has at
most \(4\Cu L\) occurrences. Each pair of unit shells
\(n\le|\Re w_j|<n+1\), \(n\ge1\), has at most
\(2\Cu\Lambda(n)\), by translating their unit intervals into
the physical variable. For \(\gamma>0\) and \(k=1,2\), put
\begin{equation}
 E_k(\gamma)=\sum_{j\ge0}2^{-j\gamma}(1+(j+1)\log 2)^k .
 \label{app:dyadic}
\end{equation}
For every integer \(M_0\ge1\) and real \(X\ge1\), dyadic blocks
give
\begin{align*}
 \sum_{n\ge M_0} n^{-2u}\Lambda(n)
    &\le M_0^{1-2u}\Lambda(M_0) E_1(2u-1),\\
 \int_X^\infty x^{-2u}\Lambda(x)^k\dd x
    &\le X^{1-2u}\Lambda(X)^k E_k(2u-1).
\end{align*}
For blocks starting at \(s=2^jM_0\) or \(2^jX\), bound the
number or length of the terms by \(s\), and use
\(\Lambda(2^{j+1}x)\le\Lambda(x)+(j+1)\log2\) for \(x\ge1\).

An occurrence of \(Z\) has modulus
\(|\mu_j|^u\exp(v\Im\Log Q_j)\).
The collar gives \(|\mu_j|\le q_0/a^2<1\), and on shell \(n\ge1\)
one also has \(|\mu_j|\le q_0\ell^2/n^2\). Thus for
\(1/2<b\le u\le r<\infty\),
\begin{equation}
\begin{split}
 |Z(p,t)|&\le A_b \exp(\Theta|v|),\\
 A_b&=4\Cu L(q_0/a^2)^b+
        2\Cu(\ell^2q_0)^b E_1(2b-1)\Lambda(1).
\end{split}
\label{app:value-envelope}
\end{equation}
The right side bounds the sum of moduli. The shell bound gives
normal convergence on compact subsets of \(u>1/2\) and uniform
tails over the real tube at fixed \(T\). An omitted occurrence
when \(N>4T+8\) satisfies \(|\Re w_j|>N-1-2T>N/2\) and has a
principal log. With \(M_0=\lfloor N/2\rfloor\ge N/3\), this gives
\begin{equation}
 |Z-Z_N|\le 2\Cu(\ell^2q_0)^u E_1(2b-1)3^{2r-1}
                N^{1-2u}\Lambda(N)\exp(\pi|v|/2).
 \label{app:divisor-tail}
\end{equation}
For an individual integer \(u=n\) use \(E_1(2n-1)3^{2n-1}\) in
this bound. The same constant without \(3^{2r-1}\) works for
a centered shell cutoff with all \(|\Re w_j|<N\) retained,
provided \(N\ge8\). Here \(\Lambda(N)=O(\log N)\) when \(N>4T+8\).

\subsection{Safe cuts and removal of the remote connectors}\label{app:clearance}

In a physical unit interval \([U,U+1]\), the zero projections
in \([U-1,U+2]\) number at most \(N_U=3\Cu\log(4+|U|)\).
Deleting the intervals of radius
\(\eta_U=1/[4(1+N_U)]\) about them removes length at most \(1/2\).
The remaining interval contains a cut \(v_0\) at distance at
least \(\eta_U\) from all projected zeros, since projections
outside the enlarged interval have distance at least one.
Translating the finite cut windows at the reference time gives
the cuts of Section~\ref{sec:complex}. Compactness of the
connectors and continuity of \(H(t,c(t)+w)\) preserve clearance
in a smaller time interval. For remote closures we need only
a sequence of cuts at each slice.

To estimate at such a cut put \(x=v_0-c(t)\) and
\(\lambda_0=\log(3+|v_0|)\). The once-subtracted product
logarithmic derivative gives
\[
 h(t,x+\ii y)-h(t,x+4\ii)
  =\sum_\rho\left(\frac1{v_0+\ii y-\rho}-\frac1{v_0+4\ii-\rho}\right).
\]
The regularizers cancel before summation. For \(|y|\le4\)
the near shell \(|\Re\rho-v_0|<1\) has
\(2\Cu\lambda_0\) as a count bound. Its denominator factors
have moduli at least \(\eta_U\) and \(3\), and the numerator is
at most \(8\). On shell \(n\ge1\) both factors are at least \(n\)
and the count is at most \(2\Cu(\lambda_0+\log(n+1))\).
The sums \(\sum n^{-2}\le2\) and
\(\sum n^{-2}\log(n+1)\le3\) bound the far part by
\(32\Cu\lambda_0+48\Cu\). Since
\(\eta_U^{-1}\le4+24\Cu\lambda_0\), the upper safe-line
bound and the near term \(16\Cu\lambda_0/(3\eta_U)\) give
\[
 |h(t,x+\ii y)|\le240\Cu^2\lambda_0^2 .
\]
For \(|x|\ge6,\ |y|\le4\) the principal log has phase at most
\(\pi/2\) and \(|Q|\ge x^2/2\), because \(|D|\le1/16\).
The absolute integral of \(f_ph\dd w\) on the remote connector,
before Cauchy normalization, is therefore at most
\begin{equation}
 1920 \Cu^2(2\ell^2)^u |x|^{-2u}
                      \Lambda(x)^2\exp(\pi|v|/2).
 \label{app:edge-bound}
\end{equation}
This tends to zero for the \(u\) in Theorem~\ref{thm:complex-time}.

\subsection{Time bounds for the spatial formula}\label{app:time-bounds}

On either horizontal safe line the reciprocal-square and twice
reciprocal-cube sums for \(h_w,h_{ww}\) are absolute. The divisor
distance is at least \(3\) in the near shell and at least \(n\)
in shell \(n\ge1\). Using the same count and elementary sums
as above, including the safe-line bound for \(h\), gives
\[
 |h|,\ |h_w|,\ |h_{ww}|\le A\Lambda(x),\qquad A=24\Cu.
\]
Take \(C_c=\sup_I|c'|\) and \(C_D=\sup_I|D'|\), finite by the
target contour. The inherited deflated-field and pair bounds
give \(C_c\le16 \Cu L^2/a+190 \Cu L\) and
\(C_D\le2+56d \Cu q_0 L/a^2\) \cite{BakerGT41,BakerCM03}.
Then \(|h_t|\le C_h\Lambda(x)^2\) with
\(C_h=A(1+C_c)+2A^2\). For \(J_p=1+|p|\) and \(|x|\ge6\),
\[
 \begin{aligned}
 |f_p|&\le (2\ell^2)^u|x|^{-2u}\exp(\pi|v|/2),\\
 |(f_p)_t|&\le 2C_D J_p|f_p|/x^2,\qquad
 |(f_p)_w|\le 8J_p|f_p|/|x|,\qquad
 |(f_p)_{ww}|\le 68J_p^2|f_p|/x^2.
 \end{aligned}
\]
In fact \((f_p)_w=-2pw f_p/Q\) and
\((f_p)_{ww}=(-2p/Q+4p(p+1)w^2/Q^2)f_p\).
The displayed bounds follow using \(|w|\le2|x|\) and
\(|Q|\ge x^2/2\). Either differentiated integrand of
\eqref{eq:complex-time} is bounded on these lines by
\(C_{\rm sl}J_p^2 |f_p|\Lambda(x)^2\), where
\[
 C_{\rm sl}=C_h+A(2C_D+68+8C_c)+8A^2 .
\]
For the value integrand use \(A|f_p|\Lambda(x)\).
The spatial primitive in the proof of
Theorem~\ref{thm:complex-time} tends to zero at the infinite
endpoints by these same estimates.

On a compact subchart let \(P_0\) be the total finite contour
length, \(W_0=\max|w|\) and \(\kappa=\min|Q|>0\) there.
Let \(A_0\ge1\) bound \(|h|,|h_w|,|h_{ww}|\) on these pieces
over the subchart. Clearance makes each constant finite. Set
\[
\begin{split}
 R_{b,r}&=\max_{q=b,r}(\ell^2/\kappa)^q,\\
 C_{\rm ch}&=A_0(1+C_c+C_D/\kappa+2/\kappa+
                     4W_0^2/\kappa^2+2C_cW_0/\kappa)
                        +(2+2W_0/\kappa)A_0^2 .
\end{split}
\]
The finite contribution to the value formula and each
differentiated formula is at most
\((P_0/(2\pi))R_{b,r} C_{\rm ch}J_p^2\exp(\Theta|v|)\),
by the explicit derivatives of \(f_p,h\).
Bounding the four full rays by the integrals over \(|x|\ge6\)
then gives
\begin{equation}
\begin{split}
 |Z_t| &\le C_Z J_p^2\exp(\Theta|v|),\\
 C_Z &=\frac{P_0 R_{b,r} C_{\rm ch}}{2\pi}
   +\frac{2 C_{\rm sl}}{\pi}\max_{q=b,r}((2\ell^2)^q6^{1-2q})
                          E_2(2b-1)\Lambda(6)^2 .
\end{split}
\label{app:chart-time}
\end{equation}
The maximum arises because the scale factor is exponential in
the real part \(u\). On the fixed finite pieces the integrands are
\(C^1\) in physical time by clearance. On each safe-line point
they are continuous in time and the displayed estimates
give an integrable majorant on compact \(p\)-sets and compact
time subcharts. This proves the asserted differentiation,
continuity and holomorphy in Theorem~\ref{thm:complex-time}.
The finite chart constants concern fixed physical height.

For \(X\ge7\) let \(Z^{[X]}=\mathcal S^{[X]}(f_ph)\) with the
finite pieces retained and the four rays stopped at \(|x|=X\).
The dyadic integral estimates and \(2\ell^2<1\) give, on
\(b\le u\le r\),
\begin{equation}
\begin{aligned}
 |Z-Z^{[X]}|&\le U_0(X)\exp(\pi|v|/2),\\
 |Z_t-Z_t^{[X]}|&\le U_1(X)J_p^2\exp(\pi|v|/2),\\
 U_0(X)&=\frac{2A}{\pi}(2\ell^2)^b E_1(2b-1)
                                X^{1-2b}\Lambda(X),\\
 U_1(X)&=\frac{2C_{\rm sl}}{\pi}(2\ell^2)^b E_2(2b-1)
                                X^{1-2b}\Lambda(X)^2 .
\end{aligned}
\label{app:spatial-tails}
\end{equation}
The derivative of the cut value uses the unintegrated
derivative in \eqref{eq:complex-time}. The factor \(2/\pi\)
accounts for the four rays. The bounds hold with \(u\) in
the scale and power of \(X\) in place of \(b\), since the
dyadic constants are decreasing in \(2u-1\). This gives the
fixed-height \(X^{1-2u}\log X\) and \(X^{1-2u}\log^2 X\) rates
for \(X>4T+8\).

\subsection{Errors on the damped index contours}\label{app:cauchy-errors}

For a fixed strip in Section~\ref{sec:cauchy}, set
\[
 B_N=2\Cu(\ell^2q_0)^b E_1(2b-1)3^{2r-1}
                                           N^{1-2b}\Lambda(N).
\]
Then the divisor error is bounded by \(B_N\exp(\pi|v|/2)\)
throughout the strip. On the complete vertical difference,
\begin{equation}
 |\mathcal V_{b,r}(K(Z-Z_N))|\le
   \frac{2 C_\epsilon B_N}{\pi\lambda}
                   \sqrt{\frac{\pi}{\epsilon}}\exp(\pi^2/(16\epsilon)).
 \label{app:weighted-divisor-tail}
\end{equation}
Indeed \(\sqrt{\lambda^2+v^2}\ge\lambda\), and translating the
Gaussian integral after completing the square pays the factor
\(\exp(\pi^2/(16\epsilon))\).
This rate concerns the whole strip. Evaluating the signed residue
values gives the sharper \eqref{cau:residue-error}.

Write \(U=1+r\) so that \(J_p\le U+|v|\), and set
\(C_{\epsilon,t}=16C_D\ell^{-2}\exp(4\epsilon)\) for \(K_\epsilon\)
and \(C_{\epsilon,t}=0\) for \(G_\epsilon\). Formula
\eqref{cau:kernel-time} gives the same kernel bounds for \(K_t\)
with this constant in place of \(C_\epsilon\). The total absolute
bound for the differentiated horizontal integrals is
\[
 \frac{r-b}{\pi Y}\big[C_\epsilon C_Z(U+Y)^2+
                         C_{\epsilon,t}A_b\big]\Phi_\Theta(Y).
\]
For \(Y>\Theta/(2\epsilon)\) the differentiated vertical tail is
bounded by
\[
 \frac{2\Phi_\Theta(Y)}{\pi}\left[
 C_\epsilon C_Z(1+U/Y)^2
     \left(\frac{Y}{\kappa_\Theta(Y)}+\frac{1}{\kappa_\Theta(Y)^2}\right)
 +\frac{C_{\epsilon,t}A_b}{Y\kappa_\Theta(Y)}
 \right].
\]
Use \((U+v)^2/v\le(1+U/Y)^2v\) for \(v\ge Y\) and integrate
against \(\exp(-\kappa_\Theta(Y)(v-Y))\). This pays both the
complete derivative and the derivative of the finite rectangle
as its height grows.

Finally \eqref{app:spatial-tails} bounds the complete vertical
derivative error due to cutting the independent spatial rays by
\begin{equation}
 \frac{2}{\pi\lambda}\int_0^\infty
 \big[C_\epsilon U_1(X)(U+v)^2+C_{\epsilon,t}U_0(X)\big]
               \exp(-\epsilon v^2+\pi v/2)\dd v .
 \label{app:weighted-spatial-tail}
\end{equation}
For the value error replace the square bracket by
\(C_\epsilon U_0(X)\). All Gaussian moments are finite at the
stated fixed damping. A finite maximum of the compact chart
constants suffices for the established real tube at each fixed \(T\).
The estimates use real physical time and fixed residue indices.

